\title{Group Sequence Policy Optimization}

\begin{abstract}
We present Group Sequence Policy Optimization (GSPO), a reinforcement learning algorithm designed to be stable, efficient, and effective for training large language models. In contrast to prior approaches that utilize token-level importance ratios, GSPO introduces an importance ratio grounded in the overall sequence likelihood, implementing clipping, reward computation, and optimization at the sequence level. Our experiments show that GSPO not only surpasses GRPO in terms of training efficiency and performance but also considerably enhances the stability of Mixture-of-Experts (MoE) RL training and can simplify RL system infrastructure. These advantages of GSPO have played a pivotal role in the significant advancements seen in the latest Qwen3 models.
\end{abstract}

\section{Introduction}
\label{sec:intro}

Reinforcement learning (RL) has become a cornerstone framework for advancing the capabilities of language models at scale \citep{o1,dpsk-r1,qwq32b,qwen3}. Utilizing RL on a large scale enables language models to engage in more complex tasks, including those found in competitive mathematics and programming, by supporting more elaborate and extended reasoning sequences.

Ensuring stable and resilient training dynamics is essential for scaling RL via increased computational resources. Nonetheless, prevailing advanced RL algorithms such as GRPO \citep{grpo} suffer from pronounced stability problems when applied to large-scale language model training, frequently causing dramatic and irrecoverable model collapse \citep{qwen3, minimax-m1}. These instability challenges represent a significant obstacle to extending the performance limits of language models through prolonged RL training.

In this work, we demonstrate that GRPO's instability arises from a core flaw: the improper utilization and breakdown of importance sampling weights within its algorithmic structure. This flaw results in substantial high-variance noise during training, which intensifies as the length of the generated responses grows. The situation is further exacerbated by the incorporation of the clipping procedure, culminating in the collapse of the model.

In response to these fundamental challenges, we introduce \textbf{Group Sequence Policy Optimization (GSPO)}, a novel reinforcement learning algorithm designed for the training of large language models. GSPO’s principal contribution is its rigorously derived importance ratio, formulated according to sequence likelihood \citep{click}, which adheres to the core concept of importance sampling. Moreover, GSPO evaluates normalized rewards as advantages across several generated responses to a single query, thereby maintaining consistency between reward assignment at the sequence level and the optimization process.

Through our experiments, we observe that GSPO considerably outperforms GRPO in terms of training stability, efficiency, and overall effectiveness. Notably, GSPO intrinsically addresses the instability issues that commonly arise in the reinforcement learning (RL) training of large-scale Mixture-of-Experts (MoE) architectures, thereby removing the necessity for elaborate stabilization techniques and offering promise for streamlining RL frameworks. The advantages provided by GSPO have directly supported substantial gains in the recent Qwen3 model releases. We anticipate that GSPO will serve as a solid and scalable basis for future progress in large-scale RL-based language model training.

\section{Preliminaries}

\paragraph{Notation} 
Throughout this work, we represent an autoregressive language model with parameters $\theta$ as the policy $\pi_\theta$. Here, $x$ signifies a query and the collection of all such queries is denoted by $\mathcal{D}$. For any response $y$ corresponding to query $x$, its probability under the policy $\pi_\theta$ is given by $\pi_\theta (y | x)=\prod_{t=1}^{|y|} \pi_\theta (y_t | x, y_{<t} )$, where $|y|$ represents the length of the response $y$ in terms of tokens. Each pair $(x, y)$ consisting of a query and its response is evaluated by a verifier $r$, which assigns a reward $r(x, y)$ within the range $[0, 1]$.

\paragraph{Proximal Policy Optimization (PPO)} PPO \citep{ppo} utilizes data collected from the previous policy $\pi_{\theta_\text{old}}$ and restricts the policy improvement step to remain close to the original policy, enforced via a clipping strategy. To optimize the policy, PPO leverages the following objective, where we disregard the KL regularization component in the subsequent discussion for succinctness, given it is not central to this work:

\begin{align}
\mathcal{J}_\text{PPO}(\theta) = \mathbb{E}_{ x \sim \mathcal{D},\, y \sim \pi_{\theta_\text{old}}( \cdot | x) }
\left[ \frac{1}{|y|} \sum_{t=1}^{|y|} 
\min \left( w_{t}(\theta) \widehat{A}_{t},  \, \mathrm{clip} \left( w_{t}(\theta), 1 - {\varepsilon}, 1 + {\varepsilon}\right) \widehat{A}_{t} \right)
\right],
\end{align}

The objective $\mathcal{J}_\text{PPO}(\theta)$ calculates the expected value, where data samples $x$ are drawn from $\mathcal{D}$ and sequences $y$ are generated according to $\pi_{\theta_\text{old}}( \cdot | x)$. Within this expectation, each time step $t$ over the sequence $y$ is considered. The minimum is taken between the unaltered advantage-weighted importance ratio $w_{t}(\theta) \widehat{A}_{t}$ and its clipped version, $\mathrm{clip} \left( w_{t}(\theta), 1 - {\varepsilon}, 1 + {\varepsilon}\right) \widehat{A}_{t}$. This sum is averaged by dividing by $|y|$, the sequence length.

Here, the token $y_t$ is assigned an importance weight given by $w_{t}(\theta) = \frac{ \pi_{\theta} (y_{t} | x, y_{<t}) }{ \pi_{\theta_\text{old}} (y_{t} | x,y_{<t})}$, while the advantage $\widehat{A}_{t}$ associated with $y_t$ is approximated using a separate value estimator, and $\varepsilon$ represents the parameter controlling the range for clipping importance ratios.

A principal difficulty in applying PPO is its significant dependence on the value model. Typically, the value model is nearly as large as the policy model, leading to substantial demands on memory and computational resources. Moreover, the performance of PPO is closely tied to how accurate the value estimations are. Building a trustworthy value model is a nontrivial task; extending this reliability to accommodate longer output sequences and tasks with greater complexity proves to be an even more formidable obstacle.

\paragraph{Group Relative Policy Optimization (GRPO)} GRPO \citep{grpo} avoids reliance on a value model by evaluating the relative advantage of each candidate reply in a set corresponding to a single query. In particular, GRPO seeks to maximize the following objective:

\begin{align}
\label{equ:grpo}
\mathcal{J}_\text{GRPO}(\theta) = \mathbb{E}_{ x \sim \mathcal{D},\, \{y_i\}_{i=1}^G \sim \pi_{\theta_\text{old}}( \cdot | x) }
\left[ \frac{1}{G} \sum_{i=1}^{G} \frac{1}{|y_i|} \sum_{t=1}^{|y_i|} 
\min \left( w_{i,t}(\theta) \widehat{A}_{i,t},  \, \mathrm{clip} \left( w_{i,t}(\theta), 1 - {\varepsilon}, 1 + {\varepsilon}\right) \widehat{A}_{i,t} \right)
\right],
\end{align}

where $G$ is the number of generated responses to each query $x$ (i.e., the group size), and the importance ratio $w_{i,t}(\theta)$ and advantage $\widehat{A}_{i,t}$ of token $y_{i,t}$ are:

\begin{align}
    w_{i,t}(\theta)=\frac{ \pi_{\theta} (y_{i,t} | x, y_{i,<t}) }{ \pi_{\theta_\text{old}} (y_{i,t} | x, y_{i,<t}) },\quad\
    \widehat{A}_{i,t} = \widehat{A}_{i} = \frac{ r(x, y_i) - \mathrm{mean} \big( \{ r(x, y_i) \}_{i=1}^G \big) }{ \mathrm{std} \big( \{ r(x, y_i) \}_{i=1}^G \big) },
\end{align}

respectively, where all the tokens in $y_i$ share the same advantage as $\widehat{A}_{i}$.

\section{Motivation}
\label{sec:motivation}

As model architectures become larger, incorporate increased sparsity (such as in Mixture-of-Experts models), and produce longer output sequences, leveraging a substantial rollout batch size becomes critical for efficient hardware utilization during RL. To make training more sample-efficient, it is common to subdivide the substantial rollout dataset into several smaller mini-batches for separate gradient update steps. However, this batching strategy leads to an inherently off-policy learning environment; here, the generated responses $y$ are drawn from a previous policy, $\pi_{\theta_\text{old}}$, instead of the currently optimized policy $\pi_\theta$. As a result, mechanisms like clipping in PPO and GRPO become necessary, as they serve to exclude samples that are excessively "off-policy" from the gradient computation process.

Although techniques such as clipping seek to address the off-policy mismatch, our analysis reveals a deeper flaw inherent in GRPO: \textit{the formulation of its objective lacks well-definedness}. This inadequacy becomes especially problematic when employing large-scale models for tasks involving extended responses, frequently resulting in severe model degradation. The root of this issue lies in the improper use of importance sampling weights within the GRPO framework. In importance sampling, the intended purpose is to compute the expected value of a function $f$ under a target distribution $\pi_\text{tar}$ by reweighting data collected from a behavior distribution $\pi_\text{beh}$ as follows:

\begin{align}
\mathbb{E}_{z \sim \pi_\text{tar}} \left[ f(z) \right] = \mathbb{E}_{z \sim \pi_\text{beh}} \left[ \frac{ \pi_\text{tar} (z) }{ \pi_\text{beh} (z)} \, f(z) \right].
\end{align}

Importantly, effective correction for the mismatch between distributions via the importance weight $\frac{ \pi_\text{tar} (z) }{ \pi_\text{beh} (z)}$ depends on aggregating a large number of samples ($N \gg 1$) drawn from the behavior policy $\pi_\text{beh}$.

On the other hand, GRPO utilizes the importance weight $\frac{ \pi_{\theta} (y_{i,t} | x, y_{i,<t}) }{ \pi_{\theta_\text{old}} (y_{i,t} | x,y_{i,<t})}$ at every token location $t$. As this weighting strategy relies solely on a single drawn token $y_{i,t}$ from the subsequent-token probability $\pi_{\theta_\text{old}}( \cdot | x, y_{i, <t})$, it is ineffective for proper correction of the training distribution. Rather than fulfilling its intended purpose, this method amplifies gradient noise with high variance, especially as sequence length increases and when the clipping procedure is employed. Through our empirical investigations, we find this can provoke a destructive form of model collapse that is frequently irreversible. Once collapse has taken place, attempts to recover—such as restoring from an earlier checkpoint, painstaking adjustment of hyperparameters like clipping thresholds, increasing sequence length during generation, or varying the RL querying strategy—are typically unsuccessful.

The aforementioned observation highlights an essential flaw in the structure of GRPO. Specifically, the inadequacy of token-level importance weighting reveals a critical insight: \textbf{the granularity of the optimization target ought to align with the granularity of the reward signal}. Because rewards are assigned at the sequence level, implementing off-policy correction on individual tokens is inherently mismatched. This realization leads us to abandon objectives at the token level in favor of investigating the use of importance weighting and optimization at the \textit{sequence level} instead.

\section{Algorithm}

\subsection{GSPO: Group Sequence Policy Optimization}

Although the use of token-level importance weights $\frac{ \pi_{\theta} (y_{i,t} | x, y_{i,<t}) }{ \pi_{\theta_\text{old}} (y_{i,t} | x,y_{i,<t})}$ presents difficulties in GRPO, our analysis indicates that, within language generation tasks, the \textit{sequence-level} importance weight $\frac{ \pi_{\theta} (y | x) }{ \pi_{\theta_\text{old}} (y | x)}$ possesses a well-defined theoretical significance. Specifically, it quantifies the divergence between the likelihood of response $y$ under $\pi_{\theta_\text{old}} (\cdot | x)$ and its probability under $\pi_{\theta} (\cdot | x)$. This correspondence coheres with sequence-level rewards, providing a natural metric for reward assignment, and functions as an informative measure within the clipping process.

Based on this straightforward observation, we propose the \textbf{Group Sequence Policy Optimization (GSPO)} algorithm. GSPO employs the following sequence-level optimization objective:

\begin{align}
\mathcal{J}_\text{GSPO} (\theta) = 
\mathbb{E}_{x \sim \mathcal{D},\, \{y_i\}_{i=1}^G \sim \pi_{\theta_\text{old}}( \cdot | x)} 
\left[
\frac{1}{G} \sum_{i=1}^{G}
\min\Big( s_{i}(\theta)\, \widehat{A}_{i},\ \mathrm{clip}\big( s_{i}(\theta), 1 - \varepsilon, 1 + \varepsilon \big)\, \widehat{A}_{i} \Big)
\right],
\label{equ:gspo}
\end{align}

where we adopt the group-based advantage estimation:

\begin{align}
\widehat{A}_{i} = \frac{r(x, y_i) - \mathrm{mean} \left( \{ r(x, y_i) \}_{i=1}^G \right) }{ \mathrm{std} \left( \{ r(x, y_i) \}_{i=1}^G \right) },
\end{align}

and define the importance ratio $s_{i}(\theta)$ based on sequence likelihood \citep{click}:

\begin{align}
s_{i}(\theta) = \left( \frac{ \pi_{\theta} (y_i | x) }{ \pi_{\theta_\text{old}} (y_i | x)} \right)^{\frac{1}{|y_i|}}
=
\exp \left( \frac{1}{|y_i|} \sum_{t=1}^{|y_i|} \log \frac{ \pi_{\theta} (y_{i,t} | x, y_{i,<t}) }{ \pi_{\theta_\text{old}} (y_{i,t} | x,y_{i,<t})} \right).
\end{align}

Consequently, GSPO implements response-level clipping, as opposed to token-level clipping, to filter out highly ``off-policy'' instances during gradient computation. This approach ensures consistency with both sequence-wise reward allocation and model optimization objectives. To mitigate variance and standardize $s_{i}(\theta)$ within a consistent numerical interval, we utilize length normalization in $s_{i}(\theta)$. Without such normalization, modifications in the probability of a small subset of tokens could lead to significant volatility in the overall importance ratio at the sequence level, necessitating different clipping thresholds for responses depending on their lengths. Additionally, it is important to observe that the scale of the clipping range in GSPO generally varies—often by orders of magnitude—from that of earlier algorithms like GRPO, owing to the inherently different formulations of importance ratios in each method.

\subsection{Gradient Analysis}
\label{subsec:gradient}

We can derive the gradient of the GSPO objective as follows (clipping is omitted for brevity):

\begin{align}
\nabla_{\theta} \mathcal{J}_\text{GSPO} (\theta)
=&\ 
\nabla_{\theta} \mathbb{E}_{ x \sim \mathcal{D},\, \{y_i\}_{i=1}^G \sim \pi_{\theta_\text{old}}( \cdot | x) }
\left[ 
\frac{1}{G} \sum_{i=1}^{G} s_{i}(\theta) \widehat{A}_{i}
\right] \\
= &\ 
\mathbb{E}_{ x \sim \mathcal{D},\, \{y_i\}_{i=1}^G \sim \pi_{\theta_\text{old}}( \cdot | x) }
\left[ 
\frac{1}{G} \sum_{i=1}^{G}
s_{i}(\theta) \widehat{A}_{i}
\, \nabla_{\theta} \log s_{i}(\theta)
\right] \\
=&\ 
\mathbb{E}_{ x \sim \mathcal{D},\, \{y_i\}_{i=1}^G \sim \pi_{\theta_\text{old}}( \cdot | x) }
\left[ 
\frac{1}{G} \sum_{i=1}^{G} 
\left( \frac{ \pi_{\theta} (y_i | x) }{ \pi_{\theta_\text{old}} (y_i | x)} \right)^{\frac{1}{|y_i|}} \widehat{A}_{i} 
\frac{1}{|y_i|} \sum_{t=1}^{|y_i|} \nabla_{\theta} \log \pi_{\theta} (y_{i,t} | x, y_{i,<t}) 
\right].
\label{equ:gspo_grad}
\end{align}

For comparison, the gradient of the GRPO objective is as follows (note that $\widehat{A}_{i,t} = \widehat{A}_{i}$):

\begin{align}
\nabla_{\theta} \mathcal{J}_\text{GRPO} (\theta) 
=&\
\nabla_{\theta} \mathbb{E}_{ x \sim \mathcal{D},\, \{y_i\}_{i=1}^G \sim \pi_{\theta_\text{old}}( \cdot | x) }
\left[ \frac{1}{G} \sum_{i=1}^{G} \frac{1}{|y_i|} \sum_{t=1}^{|y_i|} 
w_{i,t}(\theta) \widehat{A}_{i,t}
\right] \\
=&\
\mathbb{E}_{ x \sim \mathcal{D},\, \{y_i\}_{i=1}^G \sim \pi_{\theta_\text{old}}( \cdot | x) }
\left[ \frac{1}{G} \sum_{i=1}^{G} \widehat{A}_{i} 
\times \frac{1}{|y_i|} \sum_{t=1}^{|y_i|} 
\frac{ \pi_{\theta} (y_{i,t} | x, y_{i,<t}) }{ \pi_{\theta_\text{old}} (y_{i,t} | x,y_{i,<t})} 
\nabla_{\theta} \log \pi_{\theta} (y_{i,t} | x, y_{i,<t})  
\right].
\end{align}

The principal difference between GSPO and GRPO centers on \textit{the manner in which they assign weights to the gradients of token log likelihoods}. Within the GRPO framework, each token is multiplied by an ``importance weight,'' namely $\frac{ \pi_{\theta} (y_{i,t} | x, y_{i,<t}) }{ \pi_{\theta_\text{old}} (y_{i,t} | x,y_{i,<t})}$. These non-uniform weights, spanning the range $(0, 1+\varepsilon]$ when $\widehat{A}_{i} > 0$, or $[1-\varepsilon, +\infty)$ when $\widehat{A}_{i} < 0$, can exert a considerable influence—especially as training continues, potentially causing unexpected side effects. On the other hand, GSPO maintains equal weighting for all tokens within a sequence, thereby removing this potential source of instability inherent to GRPO.

\subsection{GSPO-token: A Token-level Objective Variant}

In contexts such as multi-turn RL, there can be a need for more precise advantage adjustments at the token level rather than across entire sequences. To address this requirement, we present a token-specific formulation of GSPO, referred to as \textbf{GSPO-token}, which facilitates advantage assignment on a per-token basis:

\begin{align}
\mathcal{J}_\text{GSPO-token}(\theta)
=
\mathbb{E}_{ x \sim \mathcal{D},\, \{y_i\}_{i=1}^G \sim \pi_{\theta_\text{old}}( \cdot | x) }
\Bigg[ 
\frac{1}{G} \sum_{i=1}^{G} \frac{1}{|y_i|} \sum_{t=1}^{|y_i|} 
\min \bigg( s_{i,t}(\theta) \widehat{A}_{i,t},\,
\mathrm{clip} \Big( s_{i,t}(\theta), 1-{\varepsilon}, 1+{\varepsilon} \Big) \widehat{A}_{i,t} \bigg)
\Bigg],
\label{equ:gspo-token}
\end{align}

where

\begin{align}
s_{i,t}(\theta) 
= \mathrm{sg} \left[ s_{i}(\theta) \right]  \times \frac{ \pi_{\theta} (y_{i,t} | x, y_{i,<t}) }{ \mathrm{sg} \left[ \pi_{\theta} (y_{i,t} | x, y_{i,<t}) \right] },
\end{align}

and $\mathrm{sg}[\cdot]$ denotes only taking the numerical value but stopping the gradient, corresponding to the \texttt{detach} operation in PyTorch. The gradient of GSPO-token can be derived as:

\begin{align}
\nabla_{\theta} \mathcal{J}_\text{GSPO-token}(\theta)
=&\
\nabla_{\theta} \mathbb{E}_{ x \sim \mathcal{D},\, \{y_i\}_{i=1}^G \sim \pi_{\theta_\text{old}}( \cdot | x) }
\left[
\frac{1}{G} \sum_{i=1}^{G}
\frac{1}{|y_i|} \sum_{t=1}^{|y_i|}
s_{i,t}(\theta) \widehat{A}_{i,t}
\right] \\
=&\
\mathbb{E}_{ x \sim \mathcal{D},\, \{y_i\}_{i=1}^G \sim \pi_{\theta_\text{old}}( \cdot | x) }
\left[
\frac{1}{G} \sum_{i=1}^{G} s_i (\theta) \cdot \frac{1}{|y_i|} \sum_{t=1}^{|y_i|}
\widehat{A}_{i,t} \frac{ \nabla_{\theta} \pi_{\theta} (y_{i,t} | x, y_{i,<t}) }{ \pi_{\theta} (y_{i,t} | x, y_{i,<t}) }
\right] \\
=&\
\mathbb{E}_{ x \sim \mathcal{D},\, \{y_i\}_{i=1}^G \sim \pi_{\theta_\text{old}}( \cdot | x) }
\left[
\frac{1}{G} \sum_{i=1}^{G}  \left( \frac{ \pi_{\theta} (y_i | x) }{ \pi_{\theta_\text{old}} (y_i | x)} \right)^{\frac{1}{|y_i|}}
\cdot \frac{1}{|y_i|} \sum_{t=1}^{|y_i|}
\widehat{A}_{i,t} \nabla_{\theta} \log \pi_{\theta} (y_{i,t} | x, y_{i,<t})
\right].
\label{equ:gspo-token_grad}
\end{align}

Observe that the expression $\frac{ \pi_{\theta} (y_{i,t} | x, y_{i,<t}) }{ \mathrm{sg} \left[ \pi_{\theta} (y_{i,t} | x, y_{i,<t}) \right] }$ evaluates to 1 numerically; hence, $s_{i,t}(\theta)$ is, in practice, equal to $s_{i}(\theta)$. By examining Equation~\eqref{equ:gspo} alongside Equation~\eqref{equ:gspo-token}, as well as Equation~\eqref{equ:gspo_grad} with Equation~\eqref{equ:gspo-token_grad}, one finds that GSPO-token and GSPO yield equivalent results for the optimization objective, the clipping condition, and the theoretical gradient provided that all tokens within the response $y_i$ share the same advantage (specifically, $\widehat{A}_{i,t} = \widehat{A}_{i}$). Nonetheless, GSPO-token offers greater versatility since it permits individual adjustment of advantages for each token.

\section{Experiments and Discussion}

\subsection{Empirical Results}

We conduct experiments utilizing a cold-start model derived from Qwen3-30B-A3B-Base, presenting both the progression of training rewards and the associated model performance metrics on benchmarks such as AIME'24 (measured as average Pass@1 across 32 samples), LiveCodeBench (202410-202502, average Pass@1 across 8 samples), and CodeForces (evaluated by Elo Rating). For RL training, each rollout batch is divided into four mini-batches to perform gradient updates. Within GSPO, the left and right clipping thresholds specified in Equation~\eqref{equ:gspo} are set at 3e-4 and 4e-4, respectively. We utilize GRPO as a baseline, configuring its left and right clipping values in Equation~\eqref{equ:grpo} to 0.2 and 0.27, respectively—hyperparameters we have optimized for an equitable comparison. It should be highlighted that GRPO requires implementation of the Routing Replay training strategy to ensure consistent MoE RL convergence (see further discussion in \S~\ref{subsec:routing_replay}), whereas \textit{GSPO removes the necessity for this strategy}.

As depicted in Figure~\ref{fig:results}, GSPO maintains stable progress during training. Notably, \textbf{GSPO consistently enhances performance by scaling up training compute, systematically refreshing the query set, and lengthening generation outputs}. In addition, GSPO surpasses GRPO in terms of training efficiency, attaining higher training accuracy and superior benchmark results while utilizing equivalent training compute and query consumption. Lastly, the deployment of GSPO in RL training for the recent Qwen3 models serves as robust evidence of GSPO’s effectiveness in realizing the scalability benefits of RL for large language models.

\subsection{Curious Observation on Clipping Fractions}

An important difference between GSPO and GRPO lies in their respective approaches to clipping: GSPO applies clipping at the level of entire responses, whereas GRPO operates on individual tokens. As illustrated in Figure~\ref{fig:clipping}, there is a two-order-of-magnitude gap in the proportion of clipped tokens between the two methods, and this difference remains substantial even when the clipping bounds are varied. Notably, although GSPO ends up clipping a significantly higher proportion of tokens—resulting in fewer samples contributing to training or gradient computation—it nonetheless surpasses GRPO in terms of training efficiency. This somewhat paradoxical result—that the more aggressive response-level clipping of GSPO actually enhances training efficiency—suggests that the gradient estimates at the token level produced by GRPO are subject to considerable noise and less effective at exploiting sampled data. In contrast, the sequence-based framework of GSPO yields a more consistent and informative learning signal.

\subsection{Benefit of GSPO for MoE Training}
\label{subsec:routing_replay}

\paragraph{Background}
In contrast to the RL training of fully dense models, MoE models present distinctive challenges due to their sparsely activated architecture. Notably, during experiments with the GRPO algorithm, we observed that the inherent \textit{expert-activation volatility} in MoE models can impede successful convergence during RL training. Specifically, following one or more gradient steps, the particular experts selected for generating the same output can change considerably. For instance, using the 48-layer Qwen3-30B-A3B-Base, approximately 10\% of the experts chosen for an identical rollout differ between the updated policy $\pi_{\theta}$ and the prior policy $\pi_{\theta_\text{old}}$ after each RL optimization step. This instability—especially prevalent as model depth increases—induces substantial variation in the token-level importance ratios $w_{i,t}(\theta) = \frac{ \pi_{\theta} (y_{i,t} | x, y_{i,<t}) }{ \pi_{\theta_\text{old}} (y_{i,t} | x,y_{i,<t})}$, thus diminishing their reliability, as described in \S~\ref{sec:motivation} and~\ref{subsec:gradient}, and ultimately disrupts stable RL convergence.

\paragraph{Our Previous Approach} In addressing this issue, our earlier solution involved adopting the \textbf{Routing Replay} training method. In particular, we stored the set of activated experts as determined by $\pi_{\theta_\text{old}}$, then reused these expert selection patterns when evaluating $\pi_{\theta}$ to compute importance weights as $w_{i,t}(\theta) = \frac{ \pi_{\theta} (y_{i,t} | x, y_{i,<t}) }{ \pi_{\theta_\text{old}} (y_{i,t} | x,y_{i,<t})}$. This technique guarantees that for any token $y_{i,t}$, both $\pi_{\theta} (y_{i,t} | x, y_{i,<t})$ and $\pi_{\theta_\text{old}} (y_{i,t} | x,y_{i,<t})$ utilize an identical network structure, thereby stabilizing token-level importance ratios and maintaining optimization trajectories within a consistent subnetwork across successive gradient steps. As illustrated in Figure~\ref{fig:routing_replay}, Routing Replay plays a critical role in facilitating steady convergence during GRPO training of MoE models.

\paragraph{Benefit of GSPO} While Routing Replay enables effective convergence during GRPO training for MoE models, its method of reapplying routing decisions leads to increased memory usage and additional communication costs, and it may also curtail the MoE model’s usable capacity. In comparison, as illustrated in Figure~\ref{fig:results}, GSPO removes reliance on Routing Replay, allows for direct computation of the importance ratios $s_i(\theta)$, and achieves stable optimization and typical convergence behavior. The main observation underlying this approach is that GSPO is concerned exclusively with the overall sequence likelihood, $\pi_{\theta} (y_i | x)$, rather than with each token likelihood, $\pi_{\theta} (y_{i,t} | x, y_{i,<t})$. Because the MoE architecture consistently retains its core language modeling skills, dramatic variations in sequence likelihood are avoided. Consequently, GSPO directly addresses the issue of volatile expert activation found in MoE models, eliminating the necessity for more complicated mechanisms like Routing Replay. This not only streamlines and fortifies the training regime but also permits the model to fully exploit its intrinsic capacity, free from imposed limitations.

\subsection{Benefit of GSPO for RL Infrastructure}

Due to variations in numerical precision between training engines such as Megatron and inference engines like SGLang and vLLM, standard practice involves recalculating the likelihoods of generated responses with the training engine under the previous policy $\pi_{\theta_\text{old}}$. In contrast, GSPO operates on sequence-level likelihoods rather than the more granular token-level likelihoods. This characteristic suggests that GSPO is inherently less sensitive to precision mismatches. As a result, GSPO enables direct utilization of likelihoods provided by the inference engine during optimization, eliminating the necessity of re-evaluation with the training engine. This approach is particularly advantageous in use cases involving partial rollouts, multi-turn reinforcement learning, or frameworks that decouple training from inference.

\section{Conclusion}

We introduce Group Sequence Policy Optimization (GSPO), a novel reinforcement learning method tailored for training large language models. Building upon the core concept of importance sampling, GSPO formulates importance ratios through sequence likelihoods and applies sequence-wise operations for clipping, reward assignment, and optimization. Compared to GRPO, GSPO delivers significantly improved stability, efficiency, and performance during training, and proves especially effective for large-scale RL optimization in MoE architectures—playing a crucial role in the substantial progress achieved by the latest Qwen3 models. As a scalable foundation, GSPO will enable further expansion in RL capabilities, and we anticipate this will drive profound developments in artificial intelligence.

\end{document}